% ====================================================================== Tasks 2 and 3: setup
\clearpage
\section{Setup and Metrics (Tasks 2 and 3)}
Tasks 2 and 3 use the same simulator and traces as Task 1. The points below list what is new or different.

\begin{itemize}
  \item \textbf{System:} ChampSim from the PA2 repository, single core, hashed-perceptron branch predictor.
        L1I 32\,KB 8-way, L1D 48\,KB 12-way, \textbf{L2C 512\,KB 8-way (1024 sets)}, LLC 2\,MB 16-way.
        Every cache except the L2C uses LRU. There is no prefetcher unless SPP is named.
        25M warm-up and 25M simulated instructions on Trace1--Trace3.
  \item \textbf{Baseline:} every speedup is measured against no prefetcher and LRU at the L2C (the repository
        default, \code{./build\_champsim.sh no lru}). We also show SRRIP (\code{rrip.l2c\_repl}) because I-RRIP is
        built on top of it.
  \item \textbf{Metrics:} IPC = instructions / cycles, and speedup = IPC / IPC$_{\text{baseline}}$. The MPKI values are
        the ChampSim ROI statistics \code{L2C LOAD MPKI}, \code{L2C INSTRUCTION LOAD MPKI} and \code{L2C DATA LOAD MPKI}
        (instruction + data = load). L1D accuracy and coverage are defined as in Task 1. When we average over the
        traces we use the geometric mean for IPC and speedup and the arithmetic mean for everything else.
  \item \textbf{Interface change (required by Task 2A):} \code{update\_replacement\_state} now takes a ninth argument,
        \code{uint8\_t isInstruction}. We added it to the function-pointer type and the declarations in
        \code{inc/cache.h}, to every replacement file, and to all 11 call sites. Each call site passes the
        \code{instruction} field of the packet that causes the update (the MSHR entry on a fill, the RQ/WQ/PQ entry
        on a hit). A line counts as an instruction line when \code{isInstruction} is set and the request is a
        \code{LOAD} or \code{PREFETCH}. This is the same rule \code{fill\_cache()} uses for \code{BLOCK::instruction},
        and the same rule behind the L2 instruction MPKI statistic.
  \item \textbf{Two harness fixes} (neither one changes timing):
        (1) \code{PAGE\_TABLE\_WALKER::handle\_page\_fault()} is declared \code{uint64\_t} but had no \code{return}.
        GCC\,$\ge$\,8 treats the end of such a function as unreachable, so the unmodified repository built with
        GCC\,11 crashed after a few thousand instructions. We added \code{return 0;} (no caller uses the value).
        (2) The single-core path of \code{main.cc} never called \code{l2c\_replacement\_final\_stats()}. We added
        the call so the policy can print its decision statistics.
  \item \textbf{Reproducibility:} every configuration in this report is one line in \code{scripts/experiments.sh}.
        Policy parameters are compile-time macros, and their defaults are the chosen values. For a sweep point,
        \code{scripts/build.sh} writes a two-line wrapper that \code{\#define}s the macro and \code{\#include}s the policy file.
  \item \textbf{SPP and compiler:} Task 3 uses the Task 1 prefetcher (\code{prefetcher/spp.l1d\_pref}, \code{SPP-Final})
        unchanged. Tasks 2 and 3 were compiled with GCC\,11 under WSL, while Task 1 used GCC\,7 in Docker, so the
        absolute numbers differ slightly. The baseline IPC is 0.2321 / 0.2499 / 0.3820 here and
        0.2322 / 0.2511 / 0.3821 in Task 1. The speedup of SPP is 1.0083 / 1.0135 / 1.0126 here and
        1.0098 / 1.0133 / 1.0128 in Task 1. All comparisons inside Tasks 2 and 3 use the same build, so they are
        consistent with each other.
\end{itemize}

% ---------------------------------------------------------------- shared data for the Task 2 plots
% Columns: lru srrip irrip dyn. Last row is the average (geomean for speedup, arithmetic mean for MPKI).
\pgfplotstableread{
g       lru    srrip  irrip  dyn
Trace1  1.0000 0.9996 0.9962 0.9993
Trace2  1.0000 1.0020 1.0276 1.0268
Trace3  1.0000 0.9990 1.0050 1.0048
Average 1.0000 1.0002 1.0095 1.0102
}\tReplSpd
\pgfplotstableread{
g       lru   srrip irrip dyn
Trace1  17.70 17.61 18.88 17.68
Trace2  28.64 27.27 23.25 23.33
Trace3  3.96  4.04  3.90  3.91
Average 16.77 16.31 15.34 14.97
}\tReplLoad
\pgfplotstableread{
g       lru   srrip irrip dyn
Trace1  10.37 10.55 9.25  10.54
Trace2  20.14 19.32 12.33 12.53
Trace3  1.12  1.20  0.73  0.75
Average 10.54 10.36 7.44  7.94
}\tReplInstr
\pgfplotstableread{
g       lru   srrip irrip dyn
Trace1  7.33  7.06  9.63  7.14
Trace2  8.50  7.95  10.92 10.79
Trace3  2.84  2.85  3.17  3.16
Average 6.22  5.95  7.91  7.03
}\tReplData

% \replfig{columns}{legend entries}{legend name}{bar width}{caption}{label}
\newcommand{\replfig}[6]{%
\begin{figure}[H]
\centering
\begin{subfigure}{0.49\textwidth}\centering
  \gbarplot{\tReplSpd}{#1}{cycle list name=repl, bar width=#4, ymin=0.99, ymax=1.032, ylabel={Speedup over LRU}, speedline,
    legend entries={#2}, legend columns=-1, legend to name=#3}
  \caption{Speedup over LRU}
\end{subfigure}\hfill
\begin{subfigure}{0.49\textwidth}\centering
  \gbarplot{\tReplLoad}{#1}{cycle list name=repl, bar width=#4, ymin=0, ymax=30, ylabel={L2C load MPKI}}
  \caption{L2C load MPKI}
\end{subfigure}\\[4pt]
\begin{subfigure}{0.49\textwidth}\centering
  \gbarplot{\tReplInstr}{#1}{cycle list name=repl, bar width=#4, ymin=0, ymax=22, ylabel={L2 instr.\ load MPKI}}
  \caption{L2 instruction load MPKI}
\end{subfigure}\hfill
\begin{subfigure}{0.49\textwidth}\centering
  \gbarplot{\tReplData}{#1}{cycle list name=repl, bar width=#4, ymin=0, ymax=12, ylabel={L2 data load MPKI}}
  \caption{L2 data load MPKI}
\end{subfigure}\\[3pt]
\ref*{#3}
\caption{#5}
\label{#6}
\end{figure}}
